\section{Cấu trúc dữ liệu}
\noindent\textbf{Quy ước C++17}: Chỉ số từ 0, đoạn $[l,r)$. File nguồn có sẵn header chuẩn và namespace; phần bọc lặp lại được ẩn khi in.

\Algorithm{Iterative Segment Tree}
{Segment tree trên monoid có thứ tự: gán điểm, gộp đoạn, tìm biên; bản dual cộng đoạn/truy vấn điểm}
{$\mathcal{O}(\log n)$ mỗi truy vấn/cập nhật, $\mathcal{O}(n)$ dựng}
{compactcpp}{source/cpp/data_structure/iterative_segment_tree.cpp}


\Algorithm{Lazy Segment Tree (lưu ý cài đặt)}
{Quy tắc lazy propagation, hợp thành tag (Affine, Range Set) và Tree Walk $O(\log n)$}
{}
{}{}

\begin{itemize}
    \item \textbf{Quy tắc push/pull}: Gọi \texttt{push(id, l, r)} đầu \texttt{update}/\texttt{query}; gọi \texttt{pull(id)} (\texttt{st[id] = merge(st[2*id], st[2*id+1])}) sau khi cập nhật các con.
    \item \textbf{Nhiều lazy tag}:
    \begin{itemize}
        \item \textbf{Affine ($x \mapsto a \cdot x + b$)}: Áp dụng tag mới $(c,d)$ sau $(a,b) \implies c(ax+b)+d = (ca)x + (cb+d)$. \\ Tag: \texttt{laz\_a = laz\_a * c; laz\_b = laz\_b * c + d;} \\ Node: \texttt{val = val * c + (r - l + 1) * d;}
        \item \textbf{Range Set và Range Add}: Range Set ($x \leftarrow v$) ghi đè Range Add trước đó $\implies$ đặt lại \texttt{add = 0}.
    \end{itemize}
    \item \textbf{Tree Walk ($O(\log n)$)}: Nếu con trái thỏa ($\max \ge X$ hoặc tổng $\ge S$), đẩy lazy và đi trái; nếu không, điều chỉnh mục tiêu rồi đi phải.
    \item \textbf{Giới hạn}: Cấp phát $4N$ khi đánh số từ 1 (\texttt{id << 1} và \texttt{id << 1 | 1}).
\end{itemize}



\Algorithm{Sparse Segment Tree}
{Segment tree động, chỉ cấp phát node khi cần}
{$\mathcal{O}(\log n)$ mỗi truy vấn}
{compactcpp}{source/cpp/data_structure/sparse_segment_tree.cpp}

\Algorithm{Cartesian Tree}
{Cây nhị phân có thứ tự heap theo giá trị và thứ tự BST theo chỉ số}
{$\mathcal{O}(n)$ dựng bằng monotonic stack}
{}{}

\begin{itemize}
    \item \textbf{Tính chất}: Min-heap theo giá trị ($val[\text{parent}] \le val[\text{child}]$) + duyệt in-order cho chỉ số mảng gốc ($1..n$). Gốc là giá trị nhỏ nhất. Cây con $u$ ứng với đoạn cực đại $[L_u, R_u]$ mà $A[u]$ là nhỏ nhất ($size(u) = R_u - L_u + 1$).
    \item \textbf{$\mathcal{O}(n)$ Dựng bằng stack}: Giữ nhánh phải trong stack. Với $i = 1..n$: pop $u$ có $A[u] > A[i]$; node pop cuối thành $left[i] = u$; đỉnh còn lại có $right[\text{top}] = i$. Push $i$.
    \item \textbf{Ứng dụng chính}:
    \begin{itemize}
        \item \textbf{RMQ $\iff$ LCA}: Giá trị nhỏ nhất trên đoạn $\min_{l \le k \le r} A[k] = A[\text{LCA}(l, r)]$.
        \item \textbf{Hình chữ nhật lớn nhất trong histogram}: $\max_u (A[u] \cdot size(u))$.
        \item \textbf{D\&C $\to$ Tree DP}: Chuyển D\&C trên mảng, chia tại phần tử nhỏ nhất, thành DP trên cây con.
    \end{itemize}
\end{itemize}

\Algorithm{Persistent Segment Tree}
{Persistent segment tree: cập nhật điểm, truy vấn các phiên bản cũ}
{$\mathcal{O}(\log n)$ mỗi truy vấn/cập nhật}
{compactcpp}{source/cpp/data_structure/persistent_segment_tree.cpp}

\Algorithm{Parallel Binary Search}
{Binary search song song offline cho nhiều truy vấn đơn điệu}
{$\mathcal{O}((N + Q) \log (\text{range}) \log N)$}
{compactcpp}{source/cpp/data_structure/parallel_binary_search.cpp}

\Algorithm{Wavelet Tree}
{Wavelet tree: phần tử thứ $k$, tần suất và đếm giá trị trong $[L, R]$ trên đoạn}
{$\mathcal{O}(\log (\text{alphabet}))$ mỗi truy vấn}
{compactcpp}{source/cpp/data_structure/wavelet_tree.cpp}

\Algorithm{Segment Tree Beat}
{Segment Tree Beats: chmin/chmax, cộng và truy vấn tổng trên đoạn}
{$\mathcal{O}((n+q)\log^2 n)$ khấu hao}
{compactcpp}{source/cpp/data_structure/segment_tree_beats.cpp}

\Algorithm{Kd Tree}
{Cây $k$ chiều tĩnh: điểm gần nhất và truy vấn hộp đóng song song trục}
{$O(n)$ mỗi truy vấn trong trường hợp xấu nhất; cắt tỉa không gian trên cây cân bằng}
{compactcpp}{source/cpp/data_structure/kd_tree.cpp}

\Algorithm{BIT2D}
{Fenwick 2D: tổng hình chữ nhật và cập nhật điểm}
{$\mathcal{O}(\log^2 n)$ mỗi truy vấn}
{compactcpp}{source/cpp/data_structure/fenwick_2d.cpp}

\Algorithm{Treap}
{BST ngẫu nhiên: split/merge mảng và thao tác theo khóa ẩn}
{$\mathcal{O}(\log n)$ kỳ vọng mỗi thao tác}
{compactcpp}{source/cpp/data_structure/implicit_treap.cpp}

\Algorithm{Splay Tree / Link - cut Tree}
{Link-Cut Tree: truy vấn đường đi động, link, cut và đổi gốc trên rừng}
{$\mathcal{O}(\log n)$ khấu hao mỗi thao tác}
{compactcpp}{source/cpp/data_structure/link_cut_tree.cpp}

\Algorithm{Trie}
{Trie nhị phân multiset 64-bit: thêm, xóa, đếm, XOR lớn nhất}
{$\mathcal{O}(\text{length})$ mỗi truy vấn}
{compactcpp}{source/cpp/data_structure/binary_trie.cpp}

\Algorithm{Lichao Tree}
{Li Chao Tree: thêm đường thẳng online, truy vấn min/max}
{$\mathcal{O}(\log n)$ mỗi truy vấn/thêm}
{compactcpp}{source/cpp/data_structure/li_chao.cpp}

\Algorithm{Line Container}
{Convex hull của các đường thẳng động, hệ số góc và điểm truy vấn $x$ tùy ý}
{$\mathcal{O}(\log n)$ khấu hao mỗi thao tác}
{compactcpp}{source/cpp/data_structure/line_container.cpp}

\Algorithm{Arpa Trick}
{Truy vấn min/max đoạn offline bằng DSU}
{$\mathcal{O}(n \alpha(n))$}
{compactcpp}{source/cpp/data_structure/arpa_trick.cpp}

\Algorithm{Chia căn (Square Root Decomposition)}
{Chia căn: mảng chia block và thuật toán Mo theo thứ tự Hilbert}
{$\mathcal{O}(n \sqrt{n})$}
{compactcpp}{source/cpp/data_structure/block_array_mo.cpp}

\Algorithm{Hash Map (SplitMix64)}
{Hash map khóa không dấu, hash SplitMix64 với seed do người gọi chọn}
{$\mathcal{O}(1)$ kỳ vọng mỗi lần tra cứu}
{compactcpp}{source/cpp/data_structure/hash_map.cpp}

\Algorithm{Erasable Priority Queue}
{Priority queue max/min, xóa lazy phần tử tùy ý}
{$\mathcal{O}(\log n)$ mỗi thao tác}
{compactcpp}{source/cpp/data_structure/erasable_priority_queue.cpp}

\Algorithm{Cấu trúc mở rộng (Binary Doubling)}
{Kỹ thuật gộp block theo lũy thừa 2: thêm online cho cấu trúc tĩnh (Aho, CHT)}
{$\mathcal{O}(A \log n)$ thêm khấu hao, với cấu trúc tĩnh dựng được trong $O(A)$}
{compactcpp}{source/cpp/data_structure/expandable.cpp}

\Algorithm{Interval Set}
{Phân hoạch đoạn: gán đoạn, gộp các đoạn kề nhau cùng giá trị.}
{$O((k+1)\log s)$ mỗi lần gán, $k$ đoạn bị xóa, $s$ đoạn đang lưu}
{compactcpp}{source/cpp/data_structure/interval_set.cpp}
